\section{Exploration Behavior Analysis}
\label{appendix: exploration}
We examine the exploration behavior of models in various environments, as illustrated in Table \ref{table: explore}. The ability of agents to explore plays a significant role in their performance in partially-observable environments, as diverse exploration trajectories enable agents to acquire all the necessary information.
We compare the number of locations explored by models, including rooms in BabyAI, containers in AlfWorld, and places in Jericho. This metric reflects the models' exploration capabilities. Most models are unable to explore the minimum number of locations necessary to complete the goal. \texttt{GPT-3.5-Turbo} demonstrates performance comparable to \texttt{GPT-4} in this score. Among the open-weight models, \texttt{Llama2-70b} and \texttt{CodeLlama-34b} show similar performance and both outperform \texttt{Vicuna-13b-16k}, consistent with progress rate and success rate. 
Detailed analysis on exploration behavior is not currently implemented in our \BenchName{} framework since it is feasible only for some environments.
\begin{table*}[!h]
\centering
\resizebox{\linewidth}{!}{
\begin{tabular}{lcccccc}
\toprule
          \textbf{Tasks}     & \textbf{Minimum} & \textbf{GPT-4} & \textbf{GPT-3.5-Turbo} & \textbf{Llama2-70b} & \textbf{CodeLlama-34b} & 
          \textbf{Vicuna-13b-16k} \\ \hline
Babyai - UnlocktoUnlock & 3          & 1    & 1.25   & 1          &  1.25   &   1  \\
Babyai - FindObjs5      & 3          & 2    & 3.5    & 2          &   2     &    1 \\
Babyai - Keycorridor    & 3          & 3    & 2.5    & 1.5        &   1.75     &   1.25    \\ 
Alfworld    &     3    &    5.625      &   5.125       &  1.75  & 0.125 & 1 \\
Jericho - Zork1 &   5    &  5     &   5   &  1   & 5 &1  \\
Jericho - Zork2 &   6   &   6    &   2   &   1  & 3 &3  \\
Jericho - Zork3 &   11   &   6    &   4    &  3 & 2 &  1 \\
\bottomrule
\end{tabular}
}
\caption{Comparison of the number of locations(room in babyai, containers in alfworld, and places in Jericho) explored by models. The minimum column states the least number of locations need to explore on average in order to finish the tasks.  }
\label{table: explore}
\end{table*}